\chapter{Background and Related Work}
\label{chap:background}

This chapter surveys the intellectual heritage of the thirteen reasoning breeds
implemented in OLDIA, establishes the theoretical gap that motivates the
research, and reviews the three technical pillars on which the implementation
rests: process mining with OCEL~2.0, WebAssembly as a deterministic deployment
substrate, and BLAKE3 cryptographic receipt chains.

%% ────────────────────────────────────────────────────────────────────────────
\section{The Classical AI Reasoning Systems}
\label{sec:classical-ai}

The thirteen breeds instantiated in OLDIA span five decades of AI research,
from the earliest production-rule systems of the 1960s to the autoinstinct
breeds introduced in 2026.  Each breed is characterised by a distinct
computational model, a set of mathematical invariants, and a body of
domain-specific theory.  The following subsections cover historical context,
original authorship, and the core algorithmic insight for each breed.

% ── 2.1.1 ──────────────────────────────────────────────────────────────────
\subsection{MYCIN (Shortliffe, 1976)}
\label{subsec:mycin}

MYCIN~\cite{shortliffe1976} was developed at Stanford as a clinical
decision-support system for antibiotic selection in bacterial infections.  Its
central contribution is the \emph{Certainty Factor} (CF) algebra, a
non-probabilistic calculus for combining evidence from production rules.

The combination formula distinguishes three sign cases.  When both evidence
values $a$ and $b$ are positive:
\begin{equation}
  \mathrm{CF}(a,b) = a + b - ab
  \label{eq:cf-pos}
\end{equation}
When both are negative:
\begin{equation}
  \mathrm{CF}(a,b) = a + b + ab
  \label{eq:cf-neg}
\end{equation}
When the signs differ:
\begin{equation}
  \mathrm{CF}(a,b) = \frac{a + b}{1 - \min(|a|,|b|)}
  \label{eq:cf-mixed}
\end{equation}

A strict threshold of $\mathrm{CF} > 0.2$ is required before a hypothesis
advances.  The boundary is continuous and unfakeable: in the OLDIA test suite
a CF of exactly $0.200$ is rejected while $0.201$ is accepted
(\texttt{tests/breed\_adversarial.rs}, test~\texttt{mycin\_cf\_boundary}).
Equation~\ref{eq:cf-pos} expresses bounded Dempster--Shafer-style fusion:
positive evidence accumulates sublinearly, preventing runaway certainty.

% ── 2.1.2 ──────────────────────────────────────────────────────────────────
\subsection{HEARSAY-II (Erman \& Lesser, 1980)}
\label{subsec:hearsay}

HEARSAY-II~\cite{erman1980} was built for the ARPA Speech Understanding
Research project.  Its architectural innovation is the \emph{blackboard}:
a shared, structured working memory on which independent \emph{Knowledge
Sources} (KSs) collaborate without direct coupling.

Each KS registers \emph{Knowledge Source Activation Records} (KSARs).  An
opportunistic scheduler selects the highest-rated KSAR at each cycle:
\begin{equation}
  \mathrm{rating} = \mathrm{ks\_certainty} \times \mathrm{trigger\_cf}
  \label{eq:ksar-rating}
\end{equation}

Hypothesis fusion uses a \emph{noisy-OR} model: independent supporting
hypotheses combine multiplicatively on the failure side, so joint support
$= 1 - \prod_i (1 - p_i)$.  The unfakeable scheduling invariant verified in
OLDIA is that a higher-rated KSAR always fires before a lower-rated one
regardless of declaration order (\texttt{hearsay\_opportunistic\_order}).
The implementation (\texttt{hearsay.rs}, 584~lines) uses
\texttt{f32::total\_cmp} for deterministic priority-queue ordering and
re-rates stale KSARs when blackboard state changes.

% ── 2.1.3 ──────────────────────────────────────────────────────────────────
\subsection{SOAR (Newell, Laird \& Rosenbloom, 1987)}
\label{subsec:soar}

SOAR~\cite{laird1987} is a unified cognitive architecture motivated by
Newell's call for a single computational model of mind.  Its operator
selection mechanism uses a \emph{preference algebra} with six relations:
\textsc{require}, \textsc{prohibit}, \textsc{better-than},
\textsc{worse-than}, \textsc{best}, and \textsc{indifferent}.

When preferences produce a \emph{tie impasse}---no operator is strictly
preferred---SOAR creates a subgoal to resolve the impasse.  In OLDIA the
subgoal depth is bounded to two to guarantee termination; this matches the
standard bounded-rationality assumption in the SOAR literature.  The
antisymmetry invariant---circular \textsc{worse-than} preferences produce
deterministic termination rather than infinite regress---is verified in
\texttt{soar\_antisymmetry}.  The implementation (\texttt{soar.rs},
675~lines) outputs \texttt{chunk.pref} records that expose the operator
selection trace for process-mining audit.

% ── 2.1.4 ──────────────────────────────────────────────────────────────────
\subsection{Case-Based Reasoning (Kolodner, 1993; Aamodt \& Plaza, 1994)}
\label{subsec:cbr}

Case-Based Reasoning (CBR)~\cite{kolodner1993,aamodt1994} derives solutions
by analogy with past cases.  The \emph{4R cycle} is the canonical process
model: \textbf{Retrieve}, \textbf{Reuse}, \textbf{Revise}, \textbf{Retain}.

OLDIA implements all four stages.  Retrieval uses Jaccard similarity indexed
over a \texttt{BTreeMap} for deterministic ordering.  Reuse applies
substitutional adaptation.  Revise validates the adapted solution against
domain constraints.  Retain selects high-quality cases for storage, assigning
each a BLAKE3-derived case identifier to guarantee stable cross-run identity.
The identity invariant---identical query and case produce Jaccard similarity
exactly $1.0$---is verified in \texttt{cbr\_jaccard\_identity}.

% ── 2.1.5 ──────────────────────────────────────────────────────────────────
\subsection{Prolog and Robinson Unification (Robinson, 1965; Colmerauer, 1972)}
\label{subsec:prolog}

Robinson's \emph{resolution principle}~\cite{robinson1965} established that
first-order theorem proving is reducible to a single inference rule applied to
clauses in conjunctive normal form.  Colmerauer's Prolog~\cite{colmerauer1972}
operationalised this as SLD resolution over Horn clauses, producing the first
practical logic programming language.

OLDIA implements flat-term Robinson unification with positional variables
written \texttt{?N}.  The unification kernel (\texttt{prolog8/kernel.rs},
1143~lines) supports bounded SLD resolution to prevent divergence.  The
canonical unfakeable test is the grandparent derivation:
\begin{verbatim}
grandparent(?0,?2) :- parent(?0,?1), parent(?1,?2).
\end{verbatim}
Given facts \texttt{parent(alice,bob)} and \texttt{parent(bob,carol)}, the
system must derive \texttt{grandparent(alice,carol)}.  This query is
unfakeable without shared-variable unification: pattern matching alone cannot
chain through the intermediate variable \texttt{?1}.

% ── 2.1.6 ──────────────────────────────────────────────────────────────────
\subsection{STRIPS (Fikes \& Nilsson, 1971)}
\label{subsec:strips}

STRIPS~\cite{fikes1971} introduced the \emph{add/delete list} representation
for planning operators and the \emph{frame axiom} convention: an action
changes only what its lists explicitly specify.  The planner searches forward
from an initial state using iterative deepening.

Two formal properties define STRIPS correctness.  \emph{Plan soundness}
requires that every action's preconditions are satisfied in the state at the
point of application.  \emph{Plan completeness} requires that all goal
literals are true in the final state.  OLDIA verifies both properties at
runtime: a simulated execution walks the plan step-by-step, verifying each
precondition and confirming goal satisfaction
(\texttt{strips\_plan\_soundness}).

% ── 2.1.7 ──────────────────────────────────────────────────────────────────
\subsection{GPS (Newell \& Simon, 1963)}
\label{subsec:gps}

The General Problem Solver~\cite{newell1963} introduced \emph{means-ends
analysis}: an agent selects operators by comparing the current state with a
goal state, identifying the most significant difference, and applying an
operator that reduces that difference.  A \emph{difference table} maps
difference types to applicable operators.

Operator selection in OLDIA uses fewest unmet preconditions as the heuristic.
The unfakeable invariant is \emph{difference reduction monotonicity}: each
plan step must reduce or maintain the number of unmet goal facts---a plan step
that increases the difference count is a formal defect.  GPS also requires
\emph{difference coverage}: the plan must address every unmet goal fact
(\texttt{gps\_difference\_coverage}).

% ── 2.1.8 ──────────────────────────────────────────────────────────────────
\subsection{DENDRAL (Buchanan \& Feigenbaum, 1978)}
\label{subsec:dendral}

DENDRAL~\cite{buchanan1978} was the first expert system to match domain
expert performance, deriving molecular structures from mass spectrometry data.
Its inference model is \emph{hypothesis generation with hard constraints}:
a generate-and-test loop eliminates candidates that violate forbidden
structural patterns.

The \emph{forbid} constraint is absolute in OLDIA.  A candidate matching a
forbidden pattern is eliminated regardless of its plausibility score: the
highest-scoring candidate in the enumerated set is still removed if it is
forbidden (\texttt{dendral\_forbid\_absolute}).  This is the key distinction
from soft-constraint optimisation: forbid is a monotone elimination rule, not
a penalty.

% ── 2.1.9 ──────────────────────────────────────────────────────────────────
\subsection{ELIZA (Weizenbaum, 1966)}
\label{subsec:eliza}

ELIZA~\cite{weizenbaum1966} demonstrated that a simple pattern-transformation
system could create the illusion of understanding.  Its mechanism is a
\emph{keystack}: patterns are ordered by specificity priority, and the
highest-priority matching pattern triggers a response transformation.

Three invariants are verified in OLDIA.  \emph{Specificity ordering} requires
that a higher-priority pattern fires before lower-priority alternatives.
\emph{Pronoun reflection} requires that first-person pronouns in the input are
reflected correctly in the response.  \emph{Catch-all guarantee} requires that
the system always produces a response---there is no input for which the keystack
returns empty.

% ── 2.1.10 ─────────────────────────────────────────────────────────────────
\subsection{Autoinstinct Breeds (van der Aalst lineage, 2026)}
\label{subsec:autoinstinct}

Four autoinstinct breeds extend the classical canon with implementations drawn
from the same 1960s--1970s AI lineage.

\paragraph{AutoinstinctLearning.}
Inspired by STRIPS/HACKER~\cite{winston1975}, this breed implements bitwise
heuristic planning.  The unfakeable mathematical invariant is
\emph{monotone distance reduction}: the Hamming-like distance to the goal
state must be non-increasing across plan steps.  A plan that permits
backsliding is provably incorrect.

\paragraph{AutoinstinctNeurosis.}
Based on Colby's PARRY affect machine~\cite{colby1975}, this breed models
belief-conflict processing and paranoia.  Affect state transitions are
governed by belief conflicts detected in the input; the conflict severity
determines the strength of the paranoid response.

\paragraph{AutoinstinctVision.}
Derived from Winograd's SHRDLU Blocks World perception~\cite{winograd1972},
this breed implements symbolic scene understanding.  The key data structures
are the \emph{supported-by} chain (which block rests on which) and the
\emph{clear-object} predicate (a block is clear if nothing rests on it).

\paragraph{AutoinstinctSemantics.}
Based on Schank's Conceptual Dependency theory~\cite{schank1972}, this breed
maps natural language utterances to primitive acts: \textsc{atrans}
(abstract transfer), \textsc{ptrans} (physical transfer), and \textsc{mtrans}
(mental transfer).  Semantic coverage---every input clause maps to at least one
primitive act---is the core invariant.

%% ────────────────────────────────────────────────────────────────────────────
\section{The Operational Falsifiability Problem}
\label{sec:falsifiability}

Popper's falsifiability criterion~\cite{popper1959} holds that a scientific
claim is meaningful only if it is in principle refutable by observation.
Applied to AI reasoning systems, the criterion demands that claims about
reasoning behaviour be refutable by examining what the system actually does at
runtime.

The thirteen classical systems surveyed above were theoretically falsifiable:
each had a documented algorithmic model against which behaviour could in
principle be compared.  MYCIN's certainty factor algebra was published in
mathematical detail.  SOAR's preference semantics were formally specified.
STRIPS plans were structurally verifiable.  Yet in practice these systems were
\emph{operationally opaque}: there was no infrastructure for recording,
replaying, or auditing the actual sequence of reasoning steps taken during a
run.  A clinician using MYCIN could see its conclusions but could not
independently verify that the CF combination formula had been applied
correctly, that the threshold had been honoured, or that no rule had fired
out of order.

This gap between theoretical falsifiability and operational transparency is
what we term the \emph{operational falsifiability problem}.  It has three
dimensions.  First, \emph{execution opacity}: the internal reasoning trace is
not recorded in any durable, inspectable form.  Second, \emph{model-log
divergence}: even when some logging exists, there is no systematic method for
comparing the logged behaviour against the declared process model.  Third,
\emph{cryptographic unattestedness}: there is no proof that a logged trace
corresponds to a specific input/output pair, so logs can be altered or
fabricated after the fact.

The absence of audit infrastructure was not merely a practical inconvenience;
it was an epistemological deficiency.  Without durable, cryptographically
attested, process-mined event logs, claims that a system ``behaved correctly''
reduce to assertions rather than proofs.  OLDIA addresses all three dimensions:
it records every reasoning step as an OCEL~2.0 event, validates the log
against a declared DFA lifecycle model at the run choke point, and attests the
entire output---including the event log---with a BLAKE3 receipt.

%% ────────────────────────────────────────────────────────────────────────────
\section{Process Mining and OCEL 2.0}
\label{sec:process-mining}

Process mining~\cite{vanderaalst2016} is the discipline of extracting
process knowledge from event logs recorded by information systems.  It
occupies the intersection of data mining, machine learning, and formal process
modelling.  The three core activities are \emph{process discovery} (learning a
process model from a log), \emph{conformance checking} (measuring the fit
between a log and a declared model), and \emph{enhancement} (improving an
existing model using log evidence).

Conformance checking is evaluated against four quality dimensions introduced
by van der Aalst.  \emph{Fitness} measures the fraction of log traces that
are replayable on the process model.  \emph{Precision} measures the fraction
of model behaviour that is actually observed in the log (a permissive model
that allows arbitrary behaviour scores low precision).  \emph{Generalization}
measures how well the model explains unseen traces.  \emph{Simplicity} penalises
unnecessarily complex models.  The fitness threshold of $>0.85$ is the
standard gate for a valid process model in the wasm4pm ecosystem.

\subsection{Object-Centric Event Logs (OCEL 2.0)}
\label{subsec:ocel}

Classical process mining assumes a flat, case-centric log: each event belongs
to exactly one case.  This assumption fails for AI reasoning systems, in which
a single run may simultaneously evolve multiple objects---a rule, a
hypothesis, a case, a plan, and an operator are distinct objects with
independent lifecycles that interact during a single episode.

OCEL~2.0~\cite{ghahfarokhi2021} generalises the event log to the
\emph{object-centric} setting.  An OCEL event references a set of objects
rather than a single case identifier, and each object type carries its own
lifecycle model.  The key data structures are:
\begin{itemize}
  \item \texttt{OcelEvent}: event identifier, activity name, logical timestamp,
        map from object type to object identifier list.
  \item \texttt{OcelObject}: object identifier, object type, attribute map.
  \item \texttt{OcelLog}: set of events, set of objects, object-type map.
\end{itemize}

OLDIA derives an OCEL log from every breed run via the
\texttt{derive\_ocel(breed\_id, run\_id, \&[TraceStep])} function
(\texttt{src/ocel.rs}, 632~lines).  Timestamps are \emph{logical steps}
rather than wall-clock times, eliminating non-determinism from concurrency or
scheduling jitter.

\subsection{DFA Conformance Checking}
\label{subsec:dfa}

Each of the thirteen breeds has a declared \emph{lifecycle model}: a
deterministic finite automaton (DFA) specifying the legal sequence of
reasoning phases.  For example, the MYCIN lifecycle is
$\text{init} \to \text{rule\_eval} \to \text{cf\_combine} \to \text{threshold} \to \text{conclude}$.
The CBR lifecycle enforces the 4R ordering:
$\text{retrieve} \to \text{reuse} \to \text{revise} \to \text{retain}$.

The function \texttt{validate\_ocel\_alignment(log, model)} implements a DFA
phase checker: it replays the event sequence of each object against the
breed's DFA, returning a \texttt{ConformanceResult} that includes fitness,
the set of non-conforming objects, and a phase-violation trace.
\texttt{check\_temporal\_conformance(log)} additionally verifies that logical
steps are strictly monotone---no event may carry the same or smaller step
index as a preceding event for the same object.

Both checks are wired at the \texttt{run\_breed} choke point in
\texttt{wasm.rs}: a non-conforming execution is refused at the boundary
rather than allowed to produce an attested but invalid receipt.

\subsection{Why Process Mining Is the Natural Audit Framework}
\label{subsec:pm-audit}

Process mining is not merely a useful tool for auditing AI reasoning systems;
it is the \emph{natural} framework for this task, for three reasons.

First, reasoning systems are inherently process-structured: they move through
phases (hypothesis generation, evidence combination, preference resolution,
goal reduction) in a declared order.  Process mining is designed exactly to
verify whether observed event sequences conform to declared process models.

Second, OCEL~2.0 captures the multi-object nature of AI reasoning.  A single
MYCIN run involves multiple rules, each with its own CF, and multiple
hypotheses, each with their own lifecycle.  A flat event log would collapse
this structure; OCEL preserves it.

Third, the van der Aalst fitness/precision/generalization/simplicity framework
provides quantitative metrics that translate directly into pass/fail criteria
for automated CI gates.  A breed that achieves $\text{fitness} < 0.85$ fails
the conformance gate and cannot produce an attested receipt.

%% ────────────────────────────────────────────────────────────────────────────
\section{WebAssembly as a Deployment Target}
\label{sec:wasm}

WebAssembly (WASM)~\cite{haas2017} is a stack-based virtual instruction set
designed for safe, portable, near-native-speed execution.  It was standardised
by the W3C in 2019 and is now supported natively in all major browsers and
Node.js runtimes.

\subsection{Determinism Guarantees}
\label{subsec:wasm-determinism}

WASM is \emph{deterministic by specification}: given the same module, the same
linear memory initialisation, and the same sequence of host function calls, the
execution produces bit-identical output on every platform.  This property is
essential for OLDIA.  Determinism is a merge gate in the wasm4pm project: the
\texttt{breed\_determinism.rs} test suite runs each of the thirteen breeds
twice on the same input and compares outputs with bit-exact equality.  Any
non-determinism---from unordered \texttt{HashMap} iteration, floating-point
platform variance, or unseeded random number generation---is a defect that
blocks merge.

The \texttt{to\_js\_str()} function in \texttt{utilities.rs} is the
canonical WASM boundary serialiser; \texttt{to\_js()} silently returns
\texttt{\{\}} on \texttt{wasm32} targets and is forbidden for structured
output.

\subsection{Build Targets and Toolchain}
\label{subsec:wasm-toolchain}

OLDIA uses \texttt{wasm-pack} to compile the Rust crate
\texttt{wasm4pm-cognition} to both the \texttt{nodejs} target (for CI and
server-side use) and the \texttt{browser} target (for in-browser deployment).
The one-time compilation cost is approximately 21~seconds on 2026 hardware.
Once compiled, breed execution typically completes in 2--50~microseconds,
making WASM-hosted AI reasoning competitive with native execution for
interactive applications.

The TypeScript monorepo (\texttt{packages/}) wraps the WASM boundary via
the \texttt{@wasm4pm/kernel} package, which loads the compiled module through
\texttt{WasmLoader}---a singleton that must be explicitly reset between test
runs to prevent state leakage.

%% ────────────────────────────────────────────────────────────────────────────
\section{BLAKE3 and Cryptographic Receipt Chains}
\label{sec:blake3}

BLAKE3~\cite{oconnor2021} is a cryptographic hash function that is
simultaneously faster than MD5, cryptographically secure (preimage resistance,
second-preimage resistance, collision resistance), and deterministic across
platforms.  It supports tree-based parallelism and produces 256-bit digests.

\subsection{Receipt Construction}
\label{subsec:receipt}

OLDIA constructs receipts as follows.  After a breed run completes, the
\texttt{BreedOutput} struct (which contains the reasoning result, the OCEL log,
and metadata) is serialised to canonical JSON via \texttt{serde\_json}.  The
BLAKE3 hash of this serialisation is the \texttt{output\_hash}:
\begin{equation}
  \texttt{output\_hash} = \mathrm{BLAKE3}(\mathrm{serde\_json}(\texttt{BreedOutput}))
  \label{eq:output-hash}
\end{equation}

Because \texttt{BreedOutput} contains the OCEL log, the \texttt{output\_hash}
is simultaneously a receipt of the reasoning result \emph{and} a receipt of
the process trace.  A verifier who holds the input and the receipt can
independently reproduce the hash; any divergence indicates either a different
input or a different execution.

The \texttt{run\_id} is derived from the breed identifier and the output hash:
\begin{equation}
  \texttt{run\_id} = \mathrm{BLAKE3}(\texttt{breed} \;\|\; \texttt{output\_hash})
  \label{eq:run-id}
\end{equation}

The \texttt{replay\_pointer} is the first 16 characters of the
\texttt{output\_hash}, providing a short human-readable identifier for log
correlation.

\subsection{Cryptographic Proof of Process}
\label{subsec:proof-of-process}

The receipt chain closes the operational falsifiability problem.  Theoretical
falsifiability (the algorithm is documented) plus process conformance (the
event log conforms to the declared DFA) plus cryptographic attestation (the
BLAKE3 receipt binds the log to the specific input/output pair) yields what we
term \emph{operational provability}: a third party can verify, without access
to the running system, that a specific input was processed by a specific breed
according to its declared lifecycle, and that the stated output was produced.

This is the epistemological foundation of OLDIA.  The BLAKE3 receipt does not
merely record that something happened; it constitutes evidence that a lawful
process happened.  In the terminology of the van der Aalst constitution: if the
receipt exists and the OCEL log it covers is conformant, then it worked.

%% ────────────────────────────────────────────────────────────────────────────
\section{Summary}
\label{sec:background-summary}

This chapter has established the intellectual genealogy of the thirteen OLDIA
breeds, identified the operational falsifiability problem that motivates the
research, and reviewed the three technical pillars of the implementation.

The classical AI systems from 1963 to 1993---GPS, STRIPS, MYCIN, HEARSAY-II,
SOAR, CBR, Prolog, DENDRAL, ELIZA---were theoretically falsifiable but lacked
the audit infrastructure needed for operational proof.  The four autoinstinct
breeds extend this lineage while inheriting the same gap.  OCEL~2.0 process
mining provides the natural audit framework: object-centric event logs capture
multi-object reasoning lifecycles, DFA conformance checking verifies that
declared phase orderings were honoured, and the fitness/precision metrics
provide quantitative pass/fail gates.  WebAssembly guarantees
bit-exact determinism across platforms, and BLAKE3 receipt chains
cryptographically attest that a specific input produced a specific conformant
execution.  Together these three pillars transform ``the code says it worked''
into ``the event log proves a lawful process happened.''

Chapter~\ref{chap:design} describes how these elements are combined in the
OLDIA architecture.
